\chapter*{Introduction générale}
\addcontentsline{toc}{chapter}{Introduction générale}
\markboth{Introduction générale}{}

Dans un environnement économique de plus en plus orienté vers la donnée, la
capacité d'une entreprise à analyser sa performance financière, à se situer par
rapport à ses concurrents et à anticiper ses résultats constitue un avantage
stratégique déterminant. Les états financiers, riches en information, restent
toutefois difficiles à exploiter sous leur forme brute : volumineux, hétérogènes
d'une année à l'autre, et dispersés dans des fichiers comptables peu propices à
l'analyse comparative ou prédictive.

\medskip
C'est dans ce contexte que s'inscrit le présent projet de fin d'études, réalisé
au sein de la formation \textit{Business Analytics} d'Esprit School of Business.
Il a pour objet la conception d'une \textbf{plateforme décisionnelle} dédiée à
l'analyse financière de la \textbf{Société Frigorifique et Brasserie de Tunis
(\sfbt)}, leader tunisien du secteur des boissons. L'objectif est de transformer
les états financiers de la société en un système d'aide à la décision complet,
couvrant la chaîne décisionnelle de bout en bout : de la collecte des données
jusqu'à la restitution visuelle et la prévision.

\medskip
La problématique traitée peut se résumer ainsi : \textit{comment exploiter des
états financiers hétérogènes pour produire un diagnostic financier fiable,
comparatif et prédictif de la \sfbt ?} Y répondre suppose de relever plusieurs
défis : fiabiliser des données comptables irrégulières, calculer un ensemble
cohérent de ratios financiers normalisés, comparer la \sfbt à des références
sectorielles malgré l'hétérogénéité des sources, densifier l'historique pour
permettre la modélisation, et enfin prévoir les valeurs futures.

\medskip
Pour y parvenir, la solution développée articule quatre composantes : (i) une
chaîne de traitement (\ac{ETL}) en Python assurant le nettoyage et l'unification
des données ; (ii) un moteur de calcul de \textbf{46 ratios financiers} aboutissant
à un \textbf{score global de performance} ; (iii) un \textbf{entrepôt de données}
en schéma en étoile destiné aux tableaux de bord ; et (iv) un module de
\textbf{prévision} par apprentissage automatique.

\medskip
Ce rapport est organisé en quatre chapitres. Le \textbf{premier chapitre} présente
le cadre général du projet : son contexte, sa problématique, ses objectifs et la
méthodologie adoptée. Le \textbf{deuxième chapitre} dresse un état de l'art des
concepts mobilisés : informatique décisionnelle, modélisation des entrepôts de
données, analyse financière par ratios et apprentissage automatique appliqué à la
prévision. Le \textbf{troisième chapitre} détaille l'analyse des besoins et la
conception de la solution (architecture, modèle dimensionnel, spécification des
ratios). Le \textbf{quatrième chapitre} décrit la réalisation et présente les
résultats obtenus, leur validation et leur interprétation. Une conclusion
générale synthétise les apports du travail et ouvre des perspectives
d'amélioration.
